\section*{Limitations}
\label{sec:limitations}

\paragraph{No human validation.} The most consequential gap. Every number here is produced by
an LLM: the patient, the reward, and both graders. An oracle can be perfectly repeatable and
consistently wrong, and no amount of cross-grader agreement substitutes for a certified MITI
coder rating a sample of these transcripts. We measured reproducibility (ICC(2,1)
$0.86$--$0.99$) and cross-family agreement (\S\ref{sec:heldout}) precisely because we could not
measure validity. The claim we defend is comparative --- \emph{this} arm's gains transfer to an
unrelated grader and \emph{that} arm's do not --- not that any model conducts good MI.

\paragraph{One held-out judge.} Gain retention is proposed as a diagnostic on the strength of a
single second grader. A different held-out judge could produce different retention levels. What
is more robust than the level is the \emph{comparison}: both arms are measured against the same
held-out grader, and the level offset cancels in the ratio.

\paragraph{Patient and primary grader share a model.} \primaryjudge{} both plays the patient
and computes the reward, so the policy is optimising a signal generated by its own
conversational partner. This is the coupling the held-out judge is designed to detect, and the
retention results indicate it matters. It also means our absolute scores should not be read as
estimates of quality.

\paragraph{MITI is provisional.} Dependability $0.65$, only $3.6\%$ of arm-mean variance
between arms, and sign preserved on $77.5\%$ of contrasts. Every MITI-based statement,
including the competency-threshold placement in Table~\ref{tab:thresholds}, should be treated
as descriptive. Similarly, MICI's cross-judge agreement ($r=0.20$--$0.55$) supports contrasts
but not precise rates, which is why \S\ref{sec:heldout} rests on retention instead.

\paragraph{Length is an unresolved confound.} Turn length rises $2.3$--$3.4\times$ while
conversation length falls ($28.4\rightarrow20.4$ utterances in PTO, $28.8\rightarrow25.2$ in
GRPO) and scores rise. Because the instruments score whole sessions, a shorter session
concentrated on validation is also a smaller surface on which to be marked down. We did not run
a length-controlled evaluation, so we cannot partition the gain into content and verbosity, nor
separate ``shorter because more efficient'' from ``higher-scoring because shorter''. Given the
documented length bias of LLM judges \citep{dubois2024lengthcontrolled} this bears on the
magnitude of every effect reported --- though not on the retention contrast, where both arms
lengthen and only one loses retention.

\paragraph{Single run per arm.} Each configuration was trained once. Iteration-to-iteration
variation within a run (GRPO's iteration-9 dip is visible across nearly every instrument
simultaneously) suggests run-to-run variance is not negligible, and we cannot separate it from
the method effect. Cost is the reason: the study consumed ${\approx}\$317$ in API spend, and
seed replication was not affordable.

\paragraph{Scope.} One base model (Llama-3.2-1B), one simulated patient population (96
personas), one training reward, one language, one clinical modality, ten iterations, $K=0$.
Whether the retention collapse reproduces at other scales, on other rewards, or with human
patients is open.
